% year: תשפ"ד
% semester: א
% moed: א
% number: 3
% points: 14
% categories: func-analysis, derivatives
% source: exams/מבחנים/תשפד/מועד א תשפד.pdf

\begin{question}
חשבו תחומי עלייה וירידה וקיצון מוחלט של הפונקציה
\[ f(x)=\sqrt{2x-x^2} \]
בתחום הגדרתה.
\end{question}

\begin{solution}
\textbf{תחום הגדרה.} יש לדרוש $2x-x^2=x(2-x)\ge0$, כלומר $0\le x\le 2$. תחום ההגדרה הוא הקטע הסגור $[0,2]$.

\textbf{נגזרת.} בקטע הפתוח $(0,2)$ הביטוי שבשורש חיובי, ולכן לפי כלל השרשרת
\[ f'(x)=\frac{2-2x}{2\sqrt{2x-x^2}}=\frac{1-x}{\sqrt{2x-x^2}},\qquad 0<x<2. \]
(בנקודות $x=0,2$ הפונקציה אינה גזירה – הנגזרת החד-צדדית אינסופית – אבל הן נקודות קצה.)

\textbf{סימן הנגזרת.} המכנה חיובי, ולכן סימן $f'$ הוא סימן $1-x$:
\begin{itemize}
  \item עבור $0<x<1$: $f'(x)>0$, ולכן $f$ עולה ממש בקטע $[0,1]$.
  \item עבור $1<x<2$: $f'(x)<0$, ולכן $f$ יורדת ממש בקטע $[1,2]$.
\end{itemize}
(המעבר לקטעים הסגורים מוצדק כי $f$ רציפה בהם – מסקנה ממשפט לגרנז'.)

\textbf{קיצון מוחלט.} $f$ רציפה בקטע הסגור $[0,2]$, ולכן לפי משפט ויירשטראס היא מקבלת בו מקסימום ומינימום. הנקודות החשודות הן נקודת הקריטית $x=1$ ונקודות הקצה:
\[ f(0)=0,\qquad f(1)=\sqrt{2-1}=1,\qquad f(2)=0. \]
לכן
\[ \boxed{\max_{[0,2]}f=f(1)=1,\qquad \min_{[0,2]}f=f(0)=f(2)=0.} \]
(אכן $f\ge0$ תמיד כשורש, כך שהמינימום $0$ ברור גם ישירות. בנוסף, $y=\sqrt{2x-x^2}$ שקולה ל-$(x-1)^2+y^2=1,\ y\ge0$ – חצי המעגל העליון שמרכזו $(1,0)$ ורדיוסו $1$, מה שמאשר את התוצאות.)
\end{solution}
